\section{Exact-Safe BOND-MaxSim: Design and Results}
\label{sec:results-bond}

This section addresses PRQ directly. Unlike the approximate IVF pipeline of
\cref{sec:results-ivf}, the BOND-MaxSim kernel is \emph{exact-safe}: it must 
return exactly the same top-$k$ as exhaustive MaxSim, and we reject any 
wall-clock result if the top-$k$ set differs or an ordering difference exceeds 
the declared numerical tie tolerance. Both arms in every comparison here use 
float64 products and accumulation over the same float32 embeddings. As noted 
earlier, these timings are not comparable to the float32 IVF timings.

\subsection{Kernel design}
\label{sec:bond-design}

The project-local C++17 kernel stores each document in dimension-major
(vertical) order and scans dimensions incrementally in fixed checkpoints
$\{8, 16, 32, 64, 96, 128\}$. For query token $q_i$ and document token $d_j$,
after scanning the dimension prefix $[0, t)$, , we define the partial 
product and residual norms as follows:
\begin{equation}
  p_{ij}(t) = \sum_{r < t} q_i[r]\, d_j[r], \qquad
  R^{q}_{i}(t) = \bigl\lVert q_i[t{:}] \bigr\rVert_2, \qquad
  R^{d}_{j}(t) = \bigl\lVert d_j[t{:}] \bigr\rVert_2 .
\end{equation}
The Cauchy-Schwarz inequality bounds each token similarity,
$q_i^{\top} d_j \le p_{ij}(t) + R^{q}_{i}(t)\, R^{d}_{j}(t)$, which yields a
document-level upper bound on the MaxSim score:
\begin{equation}
  U_d(t) \;=\; \sum_{i} \max_{j}
    \Bigl( p_{ij}(t) + R^{q}_{i}(t)\, R^{d}_{j}(t) \Bigr)
  \;\ge\; \operatorname{score}(q, d).
  \label{eq:bond-bound}
\end{equation}
We precompute document residual norms offline for every checkpoint, and 
we initialise the inner maxima at $-\infty$ because all token similarities 
may be negative.

Pruning uses a conservative threshold. The kernel first scores a set of 
seed documents fully (the release configuration uses the first 500 documents 
in corpus order as deterministic prefix seeds). For each fully scored document, 
it computes a rounded-down lower bound on its exact score, and takes the $k$th 
largest of these lower bounds as the threshold. This threshold can only increase 
as more documents complete. A live document is pruned only when its rounded-up 
upper bound from \cref{eq:bond-bound} is strictly below the threshold. Because 
every threshold contribution understates a true score, and because the $k$th-largest 
lower bound cannot exceed the global $k$th-largest true score, and because $U_d(t)$ 
overstates the unfinished document's true score, no document that belongs to the 
true top-$k$ can be pruned. Float64 accumulation with an explicit rounding allowance 
protects these inequalities near boundaries. Before any timing, we gated the 
implementation against a float64 NumPy oracle on randomised variable-length inputs, 
all-negative similarities, exact ties, forced-pruning workloads, and invalid inputs.

\subsection{Held-out results on raw embeddings}
\label{sec:bond-raw}

\begin{table}[tb]
  \centering
  \caption{Exact-safe BOND-MaxSim versus precision-matched compiled exhaustive
    MaxSim on 40 held-out queries (float64 products and accumulation in both
    arms; four pinned physical cores; five interleaved repetitions; one
    complete online timer). Every BOND arm returned the identical ordered
    top-10 for all 40 queries; the largest absolute score deviation across all
    arms is $1.78 \times 10^{-14}$. ``Products'' is the fraction of exhaustive
    component products actually evaluated. Latency ratios always use the
    exhaustive arm interleaved in the same run. Sources are the corresponding
    SciFact raw/PCA and NFCorpus raw release artifacts in
    \texttt{results/final/}.}
  \label{tab:bond}
  \small
  \setlength{\tabcolsep}{4.5pt}
  \begin{tabular}{@{}l S[table-format=2.4] S[table-format=1.4] S[table-format=2.2] S[table-format=2.2] S[table-format=2.2]@{}}
    \toprule
    Arm & {BOND med.\ (\si{\second})} & {Exact med.\ (\si{\second})} & {Ratio ($\times$)} & {Pruned (\%)} & {Products (\%)} \\
    \midrule
    SciFact raw, prefix-500     & 42.1947 & 3.9208 & 10.76 & 22.44 & 95.47 \\
    SciFact raw, free oracle    & 45.0854 & 3.9208 & 11.50 & 36.32 & 92.43 \\
    SciFact PCA, prefix-500     & 36.9539 & 4.4111 &  8.38 & 86.30 & 69.02 \\
    SciFact PCA, free oracle    & 37.5548 & 4.4111 &  8.51 & 97.91 & 61.53 \\
    NFCorpus raw, prefix-500    & 19.1847 & 1.9501 &  9.84 & 23.36 & 94.16 \\
    \bottomrule
  \end{tabular}
\end{table}

\Cref{tab:bond} (first row) reports the primary decision measurement. On the
40 held-out SciFact queries the BOND kernel reproduces the exact ordered
top-10 for every query, prunes 22.44\% of the 207{,}320 query-document pairs,
and yet still evaluates \textbf{95.47\%} of the component products required by
exhaustive scoring. Its median online latency is \SI{42.19}{\second} versus
\SI{3.92}{\second} (p95: \SI{43.40}{\second} versus \SI{4.10}{\second}) for
the precision-matched exhaustive arm in the same interleaved run, which gives 
a 10.76-fold slowdown. This is the main negative result of the project, and 
it is a controlled slowdown measurement, not a withdrawn-style estimate.

% \paragraph{Why the bounds do not help: late pruning.}
The per-checkpoint pruning counts explain the discrepancy between the 
visible document-pruning percentage and the negligible arithmetic saving. 
Every single one of the 46,526 SciFact prunes occurred at dimension checkpoint 96 of 128, 
with none at checkpoints 8 through 64. A document pruned at dimension 96 has already paid 
75\% of its component products, so pruning 22.44\% of documents saves only 4.53\% of the arithmetic. 
Meanwhile the kernel pays for residual norms, per-pair upper bounds, threshold maintenance, 
branching, and partial-state writes on every surviving pair. These are costs that the regular, 
SIMD-friendly exhaustive kernel does not incur. On raw normalised ColBERT embeddings, the score 
signal is spread nearly uniformly across the 128 dimensions, so the Cauchy-Schwarz residual term 
stays large until almost the entire vector has been scanned. This answers RQ1 for the raw representation: 
the bounds become useful too late to offset their runtime overhead.

\subsection{Transfer to NFCorpus}

We reran the frozen raw-order configuration unchanged on NFCorpus (Table~\ref{tab:bond}, 
last row; Figure~\ref{fig:transfer}, middle and right panels). The 
diagnosis transfers: the kernel preserves the exact ordered top-10 for all 40 held-out queries 
(maximum score deviation $7.11 \times 10^{-15}$), prunes 23.36\% of pairs, evaluates 94.16\% of 
component products, and is 9.84 times slower than its same-run exhaustive arm 
(\SI{19.18}{\second} versus \SI{1.95}{\second}  median). Pruning is 
again late, though marginally earlier than on SciFact: 
730 documents prune at checkpoint 64 and 33,210 at checkpoint 96.

\begin{figure}[tb]
  \centering
  \includegraphics[width=\textwidth]{fig8_controlled_transfer.png}
  \caption{SciFact-to-NFCorpus transfer check (40 held-out queries per
    dataset; $L=100$, $\texttt{nprobe}=8$; BOND prefix seeds $=500$; same WSL
    stack, four pinned physical cores, and online boundary). Left: IVF
    exact-ranking recovery versus rerank budget; NFCorpus saturates at 0.975
    because its candidate pool misses part of the exact top-10 for 3 of 40
    queries, whereas the SciFact pool is complete. Middle and right: raw-order
    exact-safe BOND evaluates 94--95\% of exhaustive component products and
    remains about $10\times$ slower than the same-run precision-matched
    exhaustive median on both datasets.}
  \label{fig:transfer}
\end{figure}
